% titel: Blatt 5, A1: Born-Haber-Kreisprozess (Gitterenergie)

\subsection*{Gegeben und gesucht}
$\Delta H_{\mathrm{sub}}(\mathrm{Mg})=148$, $\mathrm{IE}(\mathrm{Mg\to Mg^{2+}})=2187$,
$\Delta H_{\mathrm{vap}}(\mathrm{Br_2})=31.0$, $\Delta H_{\mathrm{dis}}(\mathrm{Br_2})=193$,
$\mathrm{EA}(\mathrm{Br})=-331$, $\Delta H_{\mathrm f}(\mathrm{MgBr_2})=-524$
(jeweils \si{\kilo\joule\per\mol}). Gesucht: Gitterenergie $U$.

\subsection*{Ansatz}
Für $\mathrm{MgBr_2}$ mit zwei $\mathrm{Br}$-Atomen lautet der Born-Haber-Zyklus:
$\mathrm{Mg(s)}\to\mathrm{Mg(g)}\to\mathrm{Mg^{2+}(g)}$ und
$\mathrm{Br_2(l)}\to 2\,\mathrm{Br(g)}\to 2\,\mathrm{Br^-(g)}$; die Teilenthalpien
addieren sich zum Kreisprozess zur Bildungsenthalpie.

\subsection*{Rechnung}
Summe der übrigen Beiträge:
\begin{align*}
  \Delta H_{\mathrm{sub}}+\mathrm{IE}+\Delta H_{\mathrm{vap}}+\Delta H_{\mathrm{dis}}+2\cdot\mathrm{EA}
  &= 148+2187+31.0+193+2(-331)\\
  &= 1897\ \si{\kilo\joule\per\mol}.
\end{align*}
Damit folgt für die Gitterenergie:
\begin{equation}
  U = \Delta H_{\mathrm f} - 1897 = -524-1897 = -2421\ \si{\kilo\joule\per\mol}.
\end{equation}

\begin{ergebnis}
Die Gitterenergie von $\mathrm{MgBr_2}$ beträgt
$U \approx -2421\ \si{\kilo\joule\per\mol}$
(beim Kristallaufbau werden rund \qty{2421}{\kilo\joule\per\mol} frei).
\end{ergebnis}

\begin{kontrolle}
Die Bildungsenthalpie ist stark negativ ($-524$), das Salz ist also stabil.
Die Summe der endothermen Schritte ($148+2187+31+193=2559$) wird durch die sehr
grosse exotherme Gitterenergie von $-2421$ und die beiden Elektronenaffinitäten
überkompensiert --- ein plausibles Bild eines energiereichen Ionenkristalls.
\end{kontrolle}

\begin{merksatz}
Born-Haber: $\Delta H_{\mathrm f}$ in einzelne messbare Schritte zerlegen und
nach der einen unbekannten Grösse (hier $U$) auflösen. Vorzeichen der
Elektronenaffinitäten und die Faktorisierung pro Atom sind die klassischen Fehlerquellen.
\end{merksatz}
